\documentclass[12pt,a4paper]{report}

% =========================================================
% PACKAGES
% =========================================================
\usepackage[a4paper,
            top=0.75in,
            bottom=0.8in,
            left=0.9in,
            right=0.9in]{geometry}

\usepackage[T1]{fontenc}
\usepackage{lmodern}
\usepackage{graphicx}
\usepackage{xcolor}
\usepackage{listings}
\usepackage{fancyhdr}
\usepackage{titlesec}
\usepackage{array}
\usepackage{booktabs}
\usepackage{longtable}
\usepackage{tabularx}
\usepackage{enumitem}
\usepackage{amsmath}
\usepackage{amssymb}
\usepackage{hyperref}
\usepackage{tikz}
\usepackage{setspace}
\usepackage{lastpage}
\usepackage{float}

% =========================================================
% COLORS
% =========================================================
\definecolor{univblue}{RGB}{20,60,110}
\definecolor{lightblue}{RGB}{235,242,250}
\definecolor{codebg}{RGB}{248,248,248}
\definecolor{codegreen}{RGB}{0,110,50}
\definecolor{codeblue}{RGB}{0,70,150}
\definecolor{codegray}{RGB}{90,90,90}

% =========================================================
% HYPERLINK SETTINGS
% =========================================================
\hypersetup{
    colorlinks=true,
    linkcolor=univblue,
    urlcolor=blue,
    citecolor=univblue
}

% =========================================================
% CODE STYLE
% =========================================================
\lstdefinestyle{CStyle}{
    language=C,
    backgroundcolor=\color{codebg},
    basicstyle=\ttfamily\small,
    keywordstyle=\color{codeblue}\bfseries,
    commentstyle=\color{codegreen},
    stringstyle=\color{red!65!black},
    numbers=left,
    numberstyle=\tiny\color{codegray},
    numbersep=8pt,
    frame=single,
    rulecolor=\color{gray!40},
    breaklines=true,
    showstringspaces=false,
    tabsize=4,
    columns=fullflexible
}

\lstset{style=CStyle}

% =========================================================
% PAGE STYLE
% =========================================================
\pagestyle{fancy}
\fancyhf{}

\fancyhead[L]{\small Operating Systems Laboratory}
\fancyhead[R]{\small Department of CSE}
\fancyfoot[L]{\small Graphic Era Deemed to be University}
\fancyfoot[C]{\small Page \thepage\ of \pageref{LastPage}}
\fancyfoot[R]{\small OS Lab Manual}

\renewcommand{\headrulewidth}{0.5pt}
\renewcommand{\footrulewidth}{0.4pt}

% =========================================================
% SECTION FORMATTING
% =========================================================
\titleformat{\chapter}
{\normalfont\huge\bfseries\color{univblue}}
{}{0pt}{}

\titleformat{\section}
{\normalfont\Large\bfseries\color{univblue}}
{\thesection}{0.7em}{}

\titleformat{\subsection}
{\normalfont\large\bfseries}
{\thesubsection}{0.6em}{}

% =========================================================
% CUSTOM COMMANDS
% =========================================================
\newcommand{\experimenttitle}[2]{
\clearpage
\begin{center}
{\Large\bfseries\color{univblue} WEEK #1}\\[0.15cm]
\rule{\textwidth}{1pt}\\[0.2cm]
{\LARGE\bfseries #2}\\[0.2cm]
\rule{\textwidth}{1pt}
\end{center}
\vspace{0.4cm}
}

\newcommand{\labfield}[1]{
\noindent\textbf{#1:}\hrulefill\\[0.35cm]
}

\newcommand{\terminal}[1]{
\begin{center}
\fbox{\parbox{0.88\textwidth}{\ttfamily #1}}
\end{center}
}

% =========================================================
% DOCUMENT
% =========================================================
\begin{document}

% =========================================================
% TITLE PAGE
% =========================================================

\begin{titlepage}

\begin{center}

\vspace*{0.5cm}

% Uncomment and change path if university logo is available
% \includegraphics[width=3.2cm]{geu-logo.png}\\[0.5cm]

{\Large\bfseries
GRAPHIC ERA DEEMED TO BE UNIVERSITY}\\[0.2cm]

{\large
Dehradun, Uttarakhand, India}\\[0.5cm]

\rule{\textwidth}{1.2pt}

\vspace{0.6cm}

{\Huge\bfseries\color{univblue}
OPERATING SYSTEMS\\
LABORATORY MANUAL}

\vspace{0.4cm}

{\Large
Practical / Laboratory Course}

\vspace{0.8cm}

\begin{tabular}{rl}
\textbf{Programme:} & B.Tech. Computer Science and Engineering \\[0.25cm]
\textbf{Department:} & Computer Science and Engineering \\[0.25cm]
\textbf{Academic Session:} & 2026--27 \\[0.25cm]
\textbf{Semester:} & \rule{4cm}{0.4pt} \\[0.25cm]
\textbf{Course Code:} & \rule{4cm}{0.4pt} \\[0.25cm]
\textbf{Course Title:} & Operating Systems Laboratory
\end{tabular}

\vfill

{\large\bfseries Prepared / Coordinated By}\\[0.25cm]

{\Large\bfseries Dr. Siddhant Thapliyal}\\
Assistant Professor\\
Department of Computer Science and Engineering

\vfill

\rule{0.8\textwidth}{0.6pt}\\
{\small Department of Computer Science and Engineering\\
Graphic Era Deemed to be University, Dehradun}

\end{center}

\end{titlepage}

% =========================================================
% STUDENT INFORMATION
% =========================================================

\chapter*{Student Information}
\addcontentsline{toc}{chapter}{Student Information}

\vspace{0.5cm}

\labfield{Name of Student}
\labfield{University Roll Number}
\labfield{Student ID}
\labfield{Programme / Branch}
\labfield{Semester}
\labfield{Section}
\labfield{Batch}
\labfield{Academic Session}

\vspace{1cm}

\begin{center}
\begin{tabular}{|p{5cm}|p{5cm}|}
\hline
\centering\textbf{Faculty Signature} &
\centering\arraybackslash\textbf{Student Signature} \\
\hline
\vspace{2cm} & \\
\hline
\end{tabular}
\end{center}

% =========================================================
% CERTIFICATE
% =========================================================

\chapter*{Certificate}
\addcontentsline{toc}{chapter}{Certificate}

\vspace{1cm}

\begin{center}
\textbf{\Large CERTIFICATE}
\end{center}

\vspace{0.8cm}

\onehalfspacing

This is to certify that
\underline{\hspace{7cm}},
University Roll Number
\underline{\hspace{4cm}},
a student of
\underline{\hspace{5cm}},
has satisfactorily completed the prescribed laboratory work for the course
\textbf{Operating Systems Laboratory} during the academic session
\textbf{2026--27}.

The experiments recorded in this laboratory manual have been performed,
verified, and evaluated as per the prescribed curriculum of the Department
of Computer Science and Engineering.

\vspace{2cm}

\begin{tabularx}{\textwidth}{X X X}
\centering Student Signature &
\centering Faculty In-Charge &
\centering\arraybackslash Head of Department
\end{tabularx}

% =========================================================
% LAB OBJECTIVES
% =========================================================

\chapter*{Objectives of the Operating Systems Laboratory}
\addcontentsline{toc}{chapter}{Objectives of the Operating Systems Laboratory}

The Operating Systems Laboratory is designed to provide practical exposure
to fundamental operating-system concepts using C programming and a
Unix/Linux environment.

The major objectives are:

\begin{enumerate}
    \item To understand process creation and process management using
    \texttt{fork()}.

    \item To understand parent and child process relationships using
    \texttt{getpid()} and \texttt{getppid()}.

    \item To study process synchronization using \texttt{wait()} and related
    system calls.

    \item To understand orphan and zombie processes.

    \item To implement Inter-Process Communication (IPC) mechanisms including:
    \begin{itemize}
        \item Pipes
        \item FIFO / Named Pipes
        \item Message Queues
        \item Shared Memory
        \item Semaphores
    \end{itemize}

    \item To understand program execution using the \texttt{exec()} family of
    system calls.

    \item To implement and analyze CPU scheduling algorithms.

    \item To implement and analyze page replacement algorithms.

    \item To develop familiarity with Unix/Linux system programming.

    \item To understand how operating-system concepts are implemented through
    practical programs.
\end{enumerate}

% =========================================================
% PREREQUISITES
% =========================================================

\chapter*{Prerequisites}
\addcontentsline{toc}{chapter}{Prerequisites}

Students should have:

\begin{itemize}
    \item Basic knowledge of the C programming language.
    \item Understanding of pointers, arrays, structures and functions.
    \item Basic knowledge of process and operating-system concepts.
    \item Familiarity with Unix/Linux commands.
    \item Ability to create, edit, compile and execute files under Linux.
\end{itemize}

% =========================================================
% SOFTWARE REQUIREMENTS
% =========================================================

\chapter*{Hardware and Software Requirements}
\addcontentsline{toc}{chapter}{Hardware and Software Requirements}

\section*{Operating System}

Any Unix/Linux based operating system, such as:

\begin{itemize}
    \item Ubuntu Linux
    \item Debian
    \item Fedora
    \item Kali Linux
    \item Linux Mint
    \item Windows Subsystem for Linux (WSL)
\end{itemize}

\section*{Compiler}

GNU Compiler Collection (GCC).

\terminal{gcc --version}

\section*{Compilation}

\terminal{gcc -Wall -Wextra program.c -o program}

\section*{Execution}

\terminal{./program}

% =========================================================
% GENERAL INSTRUCTIONS
% =========================================================

\chapter*{General Laboratory Instructions}
\addcontentsline{toc}{chapter}{General Laboratory Instructions}

\begin{enumerate}
    \item Students must attend laboratory sessions regularly.

    \item Programs must be written and executed individually unless otherwise
    instructed.

    \item Students should understand the algorithm before writing the program.

    \item Predict the output before executing programs involving processes.

    \item Always check the return value of important system calls.

    \item Programs involving \texttt{fork()} must be executed carefully to
    prevent uncontrolled process creation.

    \item Students must maintain this laboratory manual regularly.

    \item Observations and outputs should be recorded during the laboratory
    session.

    \item Plagiarized programs or reports may not be accepted.

    \item Students should be prepared for viva-voce questions during each
    experiment.
\end{enumerate}

% =========================================================
% TABLE OF CONTENTS
% =========================================================

\tableofcontents

\clearpage

% =========================================================
% INDEX
% =========================================================

\chapter*{Index of Experiments}
\addcontentsline{toc}{chapter}{Index of Experiments}

\begin{longtable}{|c|p{10cm}|c|}
\hline
\textbf{Week} & \textbf{Experiment} & \textbf{Signature} \\
\hline

1 &
Demonstration of \texttt{fork()} System Call
& \\[0.6cm]
\hline

2 &
Parent Process Computes Sum of Odd Numbers and Child Process Computes
Sum of Even Numbers Using \texttt{fork()}
& \\[0.6cm]
\hline

3 &
Demonstration of \texttt{wait()} System Call
& \\[0.6cm]
\hline

4 &
Implementation of Orphan Process and Zombie Process
& \\[0.6cm]
\hline

5 &
Implementation of Pipe
& \\[0.6cm]
\hline

6 &
Implementation of FIFO / Named Pipe
& \\[0.6cm]
\hline

7 &
Implementation of Message Queue
& \\[0.6cm]
\hline

8 &
Implementation of Shared Memory
& \\[0.6cm]
\hline

9 &
Implementation of First-Come, First-Served (FCFS) CPU Scheduling
& \\[0.6cm]
\hline

10 &
Implementation of Shortest Job First (SJF) CPU Scheduling
& \\[0.6cm]
\hline

11 &
Implementation of Priority Scheduling
& \\[0.6cm]
\hline

12 &
Implementation of FIFO Page Replacement Algorithm
& \\[0.6cm]
\hline

13 &
Implementation of Least Recently Used (LRU) Page Replacement Algorithm
& \\[0.6cm]
\hline

\end{longtable}

% =========================================================
% WEEK 1
% =========================================================

\experimenttitle{1}{Demonstration of \texttt{fork()} System Call}

\section*{Aim}

To study and demonstrate process creation using the \texttt{fork()} system call
in Unix/Linux.

\section*{Learning Outcomes}

After completing this experiment, the student will be able to:

\begin{itemize}
    \item Explain process creation using \texttt{fork()}.
    \item Differentiate parent and child processes.
    \item Interpret the return value of \texttt{fork()}.
    \item Determine the number of processes created by multiple
    \texttt{fork()} calls.
\end{itemize}

\section*{Theory}

An existing Unix/Linux process can create another process using the
\texttt{fork()} system call. The newly created process is called the
\textbf{child process}, while the process that invokes \texttt{fork()} is
called the \textbf{parent process}.

The system call is invoked once but returns independently in two processes.

\begin{center}
\begin{tabular}{|c|p{8cm}|}
\hline
\textbf{Return Value} & \textbf{Meaning} \\
\hline
$<0$ & Child process creation failed. \\
\hline
$0$ & Code is executing in the child process. \\
\hline
$>0$ & Code is executing in the parent process; the returned value is the
PID of the child. \\
\hline
\end{tabular}
\end{center}

Both parent and child continue execution from the statement following
\texttt{fork()}.

\section*{System Calls Used}

\subsection*{\texttt{fork()}}

\begin{verbatim}
pid_t fork(void);
\end{verbatim}

Creates a child process.

\subsection*{\texttt{getpid()}}

\begin{verbatim}
pid_t getpid(void);
\end{verbatim}

Returns the PID of the calling process.

\subsection*{\texttt{getppid()}}

\begin{verbatim}
pid_t getppid(void);
\end{verbatim}

Returns the PID of the parent of the calling process.

% ---------------------------------------------------------

\section*{Program 1: Single Fork}

\begin{lstlisting}
#include <stdio.h>
#include <stdlib.h>
#include <unistd.h>

int main(void)
{
    pid_t pid;

    pid = fork();

    if (pid < 0)
    {
        perror("fork failed");
        return EXIT_FAILURE;
    }

    printf("LINUX\n");

    return EXIT_SUCCESS;
}
\end{lstlisting}

\section*{Expected Output}

\begin{verbatim}
LINUX
LINUX
\end{verbatim}

\section*{Explanation}

Before \texttt{fork()}, only one process exists.

After one successful \texttt{fork()}:

\[
2^1 = 2
\]

processes exist. Therefore, both processes execute the statement:

\begin{verbatim}
printf("LINUX\n");
\end{verbatim}

% ---------------------------------------------------------

\section*{Program 2: Identifying Parent and Child}

\begin{lstlisting}
#include <stdio.h>
#include <stdlib.h>
#include <unistd.h>

int main(void)
{
    pid_t pid;

    pid = fork();

    if (pid < 0)
    {
        perror("fork failed");
        return EXIT_FAILURE;
    }
    else if (pid == 0)
    {
        printf("Child Process\n");
        printf("PID  = %d\n", (int)getpid());
        printf("PPID = %d\n", (int)getppid());
    }
    else
    {
        printf("Parent Process\n");
        printf("Parent PID = %d\n", (int)getpid());
        printf("Child PID  = %d\n", (int)pid);
    }

    return EXIT_SUCCESS;
}
\end{lstlisting}

\section*{Program 3: Multiple Forks}

\begin{lstlisting}
#include <stdio.h>
#include <unistd.h>

int main(void)
{
    fork();
    printf("LINUX\n");

    fork();
    printf("UNIX\n");

    fork();
    printf("RED HAT\n");

    return 0;
}
\end{lstlisting}

\section*{Observation}

\begin{center}
\begin{tabular}{|c|c|c|}
\hline
\textbf{Stage} &
\textbf{Processes} &
\textbf{Output Count} \\
\hline
After first fork & 2 & LINUX = 2 \\
\hline
After second fork & 4 & UNIX = 4 \\
\hline
After third fork & 8 & RED HAT = 8 \\
\hline
\end{tabular}
\end{center}

For $n$ independent successive \texttt{fork()} calls:

\[
\boxed{\text{Total Processes}=2^n}
\]

and

\[
\boxed{\text{New Child Processes}=2^n-1}
\]

\section*{Result}

Process creation using single and multiple \texttt{fork()} system calls was
successfully demonstrated.

\section*{Classwork}

\begin{enumerate}
    \item Write a program that prints the PID and PPID of parent and child.
    \item Execute two successive \texttt{fork()} calls and draw the process tree.
    \item Modify the program so that the parent and child print different
    messages.
    \item Print the return value of \texttt{fork()} in both processes.
\end{enumerate}

\section*{Homework}

\begin{enumerate}
    \item Determine the number of processes created by four successive
    \texttt{fork()} calls.
    \item Write a program where the parent prints even numbers and the child
    prints odd numbers.
    \item Draw the process tree for three successive \texttt{fork()} calls.
    \item Explain why output order may differ between executions.
\end{enumerate}

\section*{Viva Questions}

\begin{enumerate}
    \item What is \texttt{fork()}?
    \item Why does \texttt{fork()} return twice?
    \item What does \texttt{fork()} return to the child?
    \item What does it return to the parent?
    \item What is \texttt{getpid()}?
    \item What is \texttt{getppid()}?
    \item Why can parent and child execute in different orders?
    \item What is Copy-on-Write?
\end{enumerate}

% =========================================================
% WEEK 2
% =========================================================

\experimenttitle{2}
{Parent Computes Sum of Odd Numbers and Child Computes Sum of Even Numbers}

\section*{Aim}

To demonstrate independent execution of parent and child processes using
\texttt{fork()}.

\section*{Theory}

After \texttt{fork()}, parent and child execute independently in separate
address spaces. Therefore, they can perform different tasks based on the return
value of \texttt{fork()}.

\section*{Program}

\begin{lstlisting}
#include <stdio.h>
#include <stdlib.h>
#include <unistd.h>
#include <sys/wait.h>

#define MAX 100

int main(void)
{
    int a[MAX];
    int n, i;
    pid_t pid;

    printf("Enter number of elements: ");
    scanf("%d", &n);

    if (n <= 0 || n > MAX)
    {
        printf("Invalid array size.\n");
        return EXIT_FAILURE;
    }

    printf("Enter %d elements: ", n);

    for (i = 0; i < n; i++)
        scanf("%d", &a[i]);

    pid = fork();

    if (pid < 0)
    {
        perror("fork failed");
        return EXIT_FAILURE;
    }
    else if (pid == 0)
    {
        int evenSum = 0;

        for (i = 0; i < n; i++)
        {
            if (a[i] % 2 == 0)
                evenSum += a[i];
        }

        printf("Child: Sum of even numbers = %d\n",
               evenSum);

        exit(EXIT_SUCCESS);
    }
    else
    {
        int oddSum = 0;

        wait(NULL);

        for (i = 0; i < n; i++)
        {
            if (a[i] % 2 != 0)
                oddSum += a[i];
        }

        printf("Parent: Sum of odd numbers = %d\n",
               oddSum);
    }

    return EXIT_SUCCESS;
}
\end{lstlisting}

\section*{Sample Output}

\begin{verbatim}
Enter number of elements: 6
Enter 6 elements: 1 2 3 4 5 6
Child: Sum of even numbers = 12
Parent: Sum of odd numbers = 9
\end{verbatim}

\section*{Result}

Parent and child processes successfully performed two different computations.

\section*{Classwork}

Modify the program so that the child calculates the sum and the parent
calculates the product of array elements.

\section*{Homework}

Write a program where the parent finds the maximum element and the child finds
the minimum element of an array.

\section*{Viva Questions}

\begin{enumerate}
    \item Do parent and child share ordinary local variables after
    \texttt{fork()}?
    \item Why is \texttt{wait()} used?
    \item What happens if \texttt{wait()} is removed?
\end{enumerate}

% =========================================================
% WEEK 3
% =========================================================

\experimenttitle{3}{Demonstration of \texttt{wait()} System Call}

\section*{Aim}

To study process synchronization using the \texttt{wait()} system call.

\section*{Theory}

A parent process can use \texttt{wait()} to suspend its execution until one of
its child processes terminates.

\begin{verbatim}
pid_t wait(int *status);
\end{verbatim}

If the child has already terminated, \texttt{wait()} returns immediately and
collects the child's termination status.

\section*{Program}

\begin{lstlisting}
#include <stdio.h>
#include <stdlib.h>
#include <unistd.h>
#include <sys/wait.h>

int main(void)
{
    pid_t pid;
    int status;

    pid = fork();

    if (pid < 0)
    {
        perror("fork failed");
        return EXIT_FAILURE;
    }
    else if (pid == 0)
    {
        printf("I am Child\n");
        printf("Child PID = %d\n", (int)getpid());

        exit(5);
    }
    else
    {
        wait(&status);

        printf("I am Parent\n");
        printf("Child PID = %d\n", (int)pid);

        if (WIFEXITED(status))
        {
            printf("Child exit status = %d\n",
                   WEXITSTATUS(status));
        }
    }

    return EXIT_SUCCESS;
}
\end{lstlisting}

\section*{Result}

The \texttt{wait()} system call was successfully used to synchronize the parent
and child processes.

\section*{Classwork}

Remove \texttt{wait()} and execute the program five times. Record the order of
the output.

\section*{Homework}

Use \texttt{waitpid()} instead of \texttt{wait()}.

% =========================================================
% WEEK 4
% =========================================================

\experimenttitle{4}{Implementation of Orphan and Zombie Processes}

\section*{Part A: Orphan Process}

\subsection*{Aim}

To demonstrate the creation of an orphan process.

\subsection*{Theory}

A child process becomes an orphan when its original parent terminates before
the child completes execution. The operating system subsequently re-parents
the child.

\subsection*{Program}

\begin{lstlisting}
#include <stdio.h>
#include <stdlib.h>
#include <unistd.h>

int main(void)
{
    pid_t pid = fork();

    if (pid < 0)
    {
        perror("fork failed");
        return EXIT_FAILURE;
    }

    if (pid == 0)
    {
        printf("Child before parent termination\n");
        printf("PID = %d, PPID = %d\n",
               (int)getpid(),
               (int)getppid());

        sleep(5);

        printf("Child after parent termination\n");
        printf("PID = %d, PPID = %d\n",
               (int)getpid(),
               (int)getppid());
    }
    else
    {
        printf("Parent PID = %d\n", (int)getpid());
        printf("Parent is terminating.\n");
        exit(EXIT_SUCCESS);
    }

    return EXIT_SUCCESS;
}
\end{lstlisting}

\section*{Part B: Zombie Process}

\subsection*{Theory}

A zombie process is a child that has completed execution but still has an entry
in the process table because its parent has not collected its termination
status.

\subsection*{Program}

\begin{lstlisting}
#include <stdio.h>
#include <stdlib.h>
#include <unistd.h>

int main(void)
{
    pid_t pid = fork();

    if (pid < 0)
    {
        perror("fork failed");
        return EXIT_FAILURE;
    }

    if (pid == 0)
    {
        printf("Child terminating...\n");
        exit(EXIT_SUCCESS);
    }
    else
    {
        printf("Parent sleeping for 20 seconds...\n");
        sleep(20);
    }

    return EXIT_SUCCESS;
}
\end{lstlisting}

During the parent's sleep, execute:

\terminal{ps -el | grep Z}

\section*{Result}

Orphan and zombie process behavior was successfully observed.

\section*{Viva Questions}

\begin{enumerate}
    \item What is an orphan process?
    \item What is a zombie process?
    \item Is a zombie process still executing?
    \item How can a zombie process be removed?
    \item What is the role of \texttt{wait()}?
\end{enumerate}

% =========================================================
% WEEK 5
% =========================================================

\experimenttitle{5}{Implementation of Pipe}

\section*{Aim}

To implement Inter-Process Communication using an unnamed pipe.

\section*{Theory}

A pipe provides a unidirectional communication channel.

\begin{verbatim}
int pipe(int fd[2]);
\end{verbatim}

where:

\begin{itemize}
    \item \texttt{fd[0]} is the read end.
    \item \texttt{fd[1]} is the write end.
\end{itemize}

\section*{Program}

\begin{lstlisting}
#include <stdio.h>
#include <stdlib.h>
#include <string.h>
#include <unistd.h>
#include <sys/wait.h>

int main(void)
{
    int fd[2];
    pid_t pid;
    char message[] = "Graphic Era";
    char buffer[100];

    if (pipe(fd) == -1)
    {
        perror("pipe");
        return EXIT_FAILURE;
    }

    pid = fork();

    if (pid < 0)
    {
        perror("fork");
        return EXIT_FAILURE;
    }

    if (pid == 0)
    {
        close(fd[0]);

        write(fd[1],
              message,
              strlen(message) + 1);

        printf("Child wrote: %s\n", message);

        close(fd[1]);
        exit(EXIT_SUCCESS);
    }
    else
    {
        close(fd[1]);

        read(fd[0],
             buffer,
             sizeof(buffer));

        printf("Parent read: %s\n", buffer);

        close(fd[0]);
        wait(NULL);
    }

    return EXIT_SUCCESS;
}
\end{lstlisting}

\section*{Result}

Communication between parent and child using a pipe was successfully
implemented.

% =========================================================
% WEEK 6
% =========================================================

\experimenttitle{6}{Implementation of FIFO / Named Pipe}

\section*{Aim}

To implement communication between processes using FIFO.

\section*{Theory}

A FIFO is a named pipe. Unlike an unnamed pipe, it can be used by unrelated
processes.

It may be created using:

\begin{verbatim}
mkfifo("myfifo", 0666);
\end{verbatim}

\section*{Writer Program}

\begin{lstlisting}
#include <stdio.h>
#include <stdlib.h>
#include <string.h>
#include <sys/stat.h>
#include <fcntl.h>
#include <unistd.h>

int main(void)
{
    int fd;
    char message[100];

    mkfifo("myfifo", 0666);

    fd = open("myfifo", O_WRONLY);

    if (fd == -1)
    {
        perror("open");
        return EXIT_FAILURE;
    }

    printf("Enter message: ");
    fgets(message, sizeof(message), stdin);

    write(fd, message, strlen(message) + 1);

    close(fd);

    return EXIT_SUCCESS;
}
\end{lstlisting}

\section*{Reader Program}

\begin{lstlisting}
#include <stdio.h>
#include <stdlib.h>
#include <fcntl.h>
#include <unistd.h>

int main(void)
{
    int fd;
    char buffer[100];

    fd = open("myfifo", O_RDONLY);

    if (fd == -1)
    {
        perror("open");
        return EXIT_FAILURE;
    }

    read(fd, buffer, sizeof(buffer));

    printf("Received: %s", buffer);

    close(fd);
    unlink("myfifo");

    return EXIT_SUCCESS;
}
\end{lstlisting}

\section*{Execution}

Open two terminals.

\textbf{Terminal 1:}

\terminal{./reader}

\textbf{Terminal 2:}

\terminal{./writer}

\section*{Result}

Inter-process communication using FIFO was successfully demonstrated.

% =========================================================
% WEEK 7
% =========================================================

\experimenttitle{7}{Implementation of Message Queue}

\section*{Aim}

To implement Inter-Process Communication using System V message queues.

\section*{Important Functions}

\begin{itemize}
    \item \texttt{ftok()} -- Generates an IPC key.
    \item \texttt{msgget()} -- Creates or accesses a message queue.
    \item \texttt{msgsnd()} -- Sends a message.
    \item \texttt{msgrcv()} -- Receives a message.
    \item \texttt{msgctl()} -- Controls or removes the queue.
\end{itemize}

\section*{Sender}

\begin{lstlisting}
#include <stdio.h>
#include <stdlib.h>
#include <string.h>
#include <sys/ipc.h>
#include <sys/msg.h>

struct message
{
    long type;
    char text[100];
};

int main(void)
{
    key_t key;
    int msgid;
    struct message msg;

    key = ftok(".", 'A');

    msgid = msgget(key, 0666 | IPC_CREAT);

    if (msgid == -1)
    {
        perror("msgget");
        return EXIT_FAILURE;
    }

    msg.type = 1;

    printf("Enter message: ");
    fgets(msg.text, sizeof(msg.text), stdin);

    msgsnd(msgid,
           &msg,
           strlen(msg.text) + 1,
           0);

    printf("Message sent.\n");

    return EXIT_SUCCESS;
}
\end{lstlisting}

\section*{Receiver}

\begin{lstlisting}
#include <stdio.h>
#include <stdlib.h>
#include <sys/ipc.h>
#include <sys/msg.h>

struct message
{
    long type;
    char text[100];
};

int main(void)
{
    key_t key;
    int msgid;
    struct message msg;

    key = ftok(".", 'A');

    msgid = msgget(key, 0666 | IPC_CREAT);

    if (msgid == -1)
    {
        perror("msgget");
        return EXIT_FAILURE;
    }

    msgrcv(msgid,
           &msg,
           sizeof(msg.text),
           1,
           0);

    printf("Received: %s", msg.text);

    msgctl(msgid, IPC_RMID, NULL);

    return EXIT_SUCCESS;
}
\end{lstlisting}

\section*{Result}

Communication using message queues was successfully implemented.

% =========================================================
% WEEK 8
% =========================================================

\experimenttitle{8}{Implementation of Shared Memory}

\section*{Aim}

To implement Inter-Process Communication using shared memory.

\section*{Theory}

Shared memory enables multiple processes to access a common memory segment.

Important system calls are:

\begin{itemize}
    \item \texttt{ftok()}
    \item \texttt{shmget()}
    \item \texttt{shmat()}
    \item \texttt{shmdt()}
    \item \texttt{shmctl()}
\end{itemize}

\section*{Writer}

\begin{lstlisting}
#include <stdio.h>
#include <stdlib.h>
#include <string.h>
#include <sys/ipc.h>
#include <sys/shm.h>

int main(void)
{
    key_t key;
    int shmid;
    char *shared;

    key = ftok(".", 'S');

    shmid = shmget(key,
                   1024,
                   0666 | IPC_CREAT);

    if (shmid == -1)
    {
        perror("shmget");
        return EXIT_FAILURE;
    }

    shared = (char *)shmat(shmid, NULL, 0);

    if (shared == (char *)-1)
    {
        perror("shmat");
        return EXIT_FAILURE;
    }

    printf("Enter data: ");
    fgets(shared, 1024, stdin);

    printf("Data written to shared memory.\n");

    shmdt(shared);

    return EXIT_SUCCESS;
}
\end{lstlisting}

\section*{Reader}

\begin{lstlisting}
#include <stdio.h>
#include <stdlib.h>
#include <sys/ipc.h>
#include <sys/shm.h>

int main(void)
{
    key_t key;
    int shmid;
    char *shared;

    key = ftok(".", 'S');

    shmid = shmget(key,
                   1024,
                   0666);

    if (shmid == -1)
    {
        perror("shmget");
        return EXIT_FAILURE;
    }

    shared = (char *)shmat(shmid, NULL, 0);

    if (shared == (char *)-1)
    {
        perror("shmat");
        return EXIT_FAILURE;
    }

    printf("Data received: %s", shared);

    shmdt(shared);

    shmctl(shmid,
           IPC_RMID,
           NULL);

    return EXIT_SUCCESS;
}
\end{lstlisting}

\section*{Result}

Shared-memory based Inter-Process Communication was successfully implemented.

% =========================================================
% WEEK 9
% =========================================================

\experimenttitle{9}{First-Come, First-Served (FCFS) CPU Scheduling}

\section*{Aim}

To implement the FCFS CPU scheduling algorithm and calculate waiting time and
turnaround time.

\section*{Theory}

FCFS is a non-preemptive scheduling algorithm in which processes are executed
in order of arrival.

\[
TAT = CT - AT
\]

\[
WT = TAT - BT
\]

where:

\begin{itemize}
    \item $TAT$ = Turnaround Time
    \item $CT$ = Completion Time
    \item $AT$ = Arrival Time
    \item $WT$ = Waiting Time
    \item $BT$ = Burst Time
\end{itemize}

\section*{Program}

\begin{lstlisting}
#include <stdio.h>

int main(void)
{
    int n, i;
    int bt[20], wt[20], tat[20];
    float avgWT = 0.0;
    float avgTAT = 0.0;

    printf("Enter number of processes: ");
    scanf("%d", &n);

    printf("Enter burst times:\n");

    for (i = 0; i < n; i++)
    {
        printf("P%d: ", i + 1);
        scanf("%d", &bt[i]);
    }

    wt[0] = 0;

    for (i = 1; i < n; i++)
        wt[i] = wt[i - 1] + bt[i - 1];

    for (i = 0; i < n; i++)
    {
        tat[i] = wt[i] + bt[i];
        avgWT += wt[i];
        avgTAT += tat[i];
    }

    printf("\nProcess\tBT\tWT\tTAT\n");

    for (i = 0; i < n; i++)
    {
        printf("P%d\t%d\t%d\t%d\n",
               i + 1,
               bt[i],
               wt[i],
               tat[i]);
    }

    printf("\nAverage Waiting Time = %.2f\n",
           avgWT / n);

    printf("Average Turnaround Time = %.2f\n",
           avgTAT / n);

    return 0;
}
\end{lstlisting}

\section*{Classwork}

Execute FCFS for:

\begin{center}
\begin{tabular}{|c|c|}
\hline
Process & Burst Time \\
\hline
P1 & 24 \\
P2 & 3 \\
P3 & 3 \\
\hline
\end{tabular}
\end{center}

\section*{Result}

FCFS scheduling was successfully implemented.

% =========================================================
% WEEK 10
% =========================================================

\experimenttitle{10}{Shortest Job First (SJF) CPU Scheduling}

\section*{Aim}

To implement non-preemptive SJF scheduling.

\section*{Theory}

SJF chooses the ready process with the shortest CPU burst.

For a fixed set of processes available simultaneously, SJF minimizes average
waiting time.

\section*{Program}

\begin{lstlisting}
#include <stdio.h>

int main(void)
{
    int n, i, j;
    int bt[20];
    int pid[20];
    int wt[20];
    int tat[20];

    float avgWT = 0.0;
    float avgTAT = 0.0;

    printf("Enter number of processes: ");
    scanf("%d", &n);

    for (i = 0; i < n; i++)
    {
        pid[i] = i + 1;

        printf("Enter burst time for P%d: ",
               i + 1);

        scanf("%d", &bt[i]);
    }

    for (i = 0; i < n - 1; i++)
    {
        for (j = i + 1; j < n; j++)
        {
            if (bt[i] > bt[j])
            {
                int temp;

                temp = bt[i];
                bt[i] = bt[j];
                bt[j] = temp;

                temp = pid[i];
                pid[i] = pid[j];
                pid[j] = temp;
            }
        }
    }

    wt[0] = 0;

    for (i = 1; i < n; i++)
        wt[i] = wt[i - 1] + bt[i - 1];

    for (i = 0; i < n; i++)
    {
        tat[i] = wt[i] + bt[i];

        avgWT += wt[i];
        avgTAT += tat[i];
    }

    printf("\nProcess\tBT\tWT\tTAT\n");

    for (i = 0; i < n; i++)
    {
        printf("P%d\t%d\t%d\t%d\n",
               pid[i],
               bt[i],
               wt[i],
               tat[i]);
    }

    printf("\nAverage WT = %.2f\n",
           avgWT / n);

    printf("Average TAT = %.2f\n",
           avgTAT / n);

    return 0;
}
\end{lstlisting}

\section*{Result}

Non-preemptive SJF scheduling was successfully implemented.

% =========================================================
% WEEK 11
% =========================================================

\experimenttitle{11}{Priority Scheduling}

\section*{Aim}

To implement non-preemptive Priority Scheduling.

\section*{Theory}

Each process is assigned a priority. The CPU is allocated to the process with
the highest priority.

In this implementation:

\[
\text{Smaller priority number} =
\text{Higher priority}
\]

\section*{Program}

\begin{lstlisting}
#include <stdio.h>

int main(void)
{
    int n, i, j;

    int pid[20];
    int bt[20];
    int priority[20];
    int wt[20];
    int tat[20];

    float avgWT = 0.0;
    float avgTAT = 0.0;

    printf("Enter number of processes: ");
    scanf("%d", &n);

    for (i = 0; i < n; i++)
    {
        pid[i] = i + 1;

        printf("P%d burst time: ", i + 1);
        scanf("%d", &bt[i]);

        printf("P%d priority: ", i + 1);
        scanf("%d", &priority[i]);
    }

    for (i = 0; i < n - 1; i++)
    {
        for (j = i + 1; j < n; j++)
        {
            if (priority[i] > priority[j])
            {
                int temp;

                temp = priority[i];
                priority[i] = priority[j];
                priority[j] = temp;

                temp = bt[i];
                bt[i] = bt[j];
                bt[j] = temp;

                temp = pid[i];
                pid[i] = pid[j];
                pid[j] = temp;
            }
        }
    }

    wt[0] = 0;

    for (i = 1; i < n; i++)
        wt[i] = wt[i - 1] + bt[i - 1];

    for (i = 0; i < n; i++)
    {
        tat[i] = wt[i] + bt[i];

        avgWT += wt[i];
        avgTAT += tat[i];
    }

    printf("\nPID\tPriority\tBT\tWT\tTAT\n");

    for (i = 0; i < n; i++)
    {
        printf("P%d\t%d\t\t%d\t%d\t%d\n",
               pid[i],
               priority[i],
               bt[i],
               wt[i],
               tat[i]);
    }

    printf("\nAverage WT = %.2f\n",
           avgWT / n);

    printf("Average TAT = %.2f\n",
           avgTAT / n);

    return 0;
}
\end{lstlisting}

\section*{Discussion}

A major problem with priority scheduling is \textbf{starvation}. A process with
low priority may wait indefinitely.

A common solution is \textbf{aging}, in which the priority of a waiting process
is gradually increased.

\section*{Result}

Priority Scheduling was successfully implemented.

% =========================================================
% WEEK 12
% =========================================================

\experimenttitle{12}{FIFO Page Replacement Algorithm}

\section*{Aim}

To implement the FIFO page replacement algorithm and calculate the number of
page faults.

\section*{Theory}

FIFO replaces the page that entered memory earliest.

The frames can conceptually be treated as a queue.

\section*{Program}

\begin{lstlisting}
#include <stdio.h>

int main(void)
{
    int pages[100];
    int frames[20];

    int n, f;
    int i, j;

    int pointer = 0;
    int faults = 0;

    printf("Enter reference string length: ");
    scanf("%d", &n);

    printf("Enter reference string:\n");

    for (i = 0; i < n; i++)
        scanf("%d", &pages[i]);

    printf("Enter number of frames: ");
    scanf("%d", &f);

    for (i = 0; i < f; i++)
        frames[i] = -1;

    printf("\nPage\tFrames\n");

    for (i = 0; i < n; i++)
    {
        int found = 0;

        for (j = 0; j < f; j++)
        {
            if (frames[j] == pages[i])
            {
                found = 1;
                break;
            }
        }

        if (!found)
        {
            frames[pointer] = pages[i];

            pointer = (pointer + 1) % f;

            faults++;
        }

        printf("%d\t", pages[i]);

        for (j = 0; j < f; j++)
        {
            if (frames[j] == -1)
                printf("- ");
            else
                printf("%d ", frames[j]);
        }

        printf("\n");
    }

    printf("\nTotal Page Faults = %d\n",
           faults);

    return 0;
}
\end{lstlisting}

\section*{Result}

FIFO page replacement was successfully implemented and page faults were
calculated.

% =========================================================
% WEEK 13
% =========================================================

\experimenttitle{13}{Least Recently Used (LRU) Page Replacement Algorithm}

\section*{Aim}

To implement the Least Recently Used page replacement algorithm.

\section*{Theory}

LRU replaces the page that has not been referenced for the longest period of
time.

The algorithm attempts to approximate the locality of reference exhibited by
programs.

\section*{Program}

\begin{lstlisting}
#include <stdio.h>

int main(void)
{
    int pages[100];
    int frames[20];
    int lastUsed[20];

    int n, f;
    int i, j;

    int faults = 0;
    int time = 0;

    printf("Enter reference string length: ");
    scanf("%d", &n);

    printf("Enter reference string:\n");

    for (i = 0; i < n; i++)
        scanf("%d", &pages[i]);

    printf("Enter number of frames: ");
    scanf("%d", &f);

    for (i = 0; i < f; i++)
    {
        frames[i] = -1;
        lastUsed[i] = -1;
    }

    printf("\nPage\tFrames\n");

    for (i = 0; i < n; i++)
    {
        int found = -1;

        time++;

        for (j = 0; j < f; j++)
        {
            if (frames[j] == pages[i])
            {
                found = j;
                break;
            }
        }

        if (found != -1)
        {
            lastUsed[found] = time;
        }
        else
        {
            int position = -1;

            for (j = 0; j < f; j++)
            {
                if (frames[j] == -1)
                {
                    position = j;
                    break;
                }
            }

            if (position == -1)
            {
                position = 0;

                for (j = 1; j < f; j++)
                {
                    if (lastUsed[j] <
                        lastUsed[position])
                    {
                        position = j;
                    }
                }
            }

            frames[position] = pages[i];
            lastUsed[position] = time;

            faults++;
        }

        printf("%d\t", pages[i]);

        for (j = 0; j < f; j++)
        {
            if (frames[j] == -1)
                printf("- ");
            else
                printf("%d ", frames[j]);
        }

        printf("\n");
    }

    printf("\nTotal Page Faults = %d\n",
           faults);

    return 0;
}
\end{lstlisting}

\section*{Result}

LRU page replacement was successfully implemented.

% =========================================================
% ADDITIONAL EXPERIMENT 1
% =========================================================

\chapter*{Additional Experiment: \texttt{fork()} and \texttt{exec()}}
\addcontentsline{toc}{chapter}
{Additional Experiment: fork() and exec()}

\section*{Aim}

To create a child using \texttt{fork()} and execute another program using
\texttt{execl()}.

\section*{Program}

\begin{lstlisting}
#include <stdio.h>
#include <stdlib.h>
#include <unistd.h>
#include <sys/wait.h>

int main(void)
{
    pid_t pid = fork();

    if (pid < 0)
    {
        perror("fork");
        return EXIT_FAILURE;
    }

    if (pid == 0)
    {
        printf("Child executing date command:\n");

        execl("/bin/date",
              "date",
              NULL);

        perror("execl");
        _exit(EXIT_FAILURE);
    }
    else
    {
        wait(NULL);

        printf("Parent executing ls -l:\n");

        execl("/bin/ls",
              "ls",
              "-l",
              NULL);

        perror("execl");
    }

    return EXIT_SUCCESS;
}
\end{lstlisting}

% =========================================================
% ADDITIONAL EXPERIMENT 2
% =========================================================

\chapter*{Additional Experiment: Implementation of \texttt{ls | wc}}
\addcontentsline{toc}{chapter}
{Additional Experiment: Implementation of ls | wc}

\section*{Aim}

To implement a Unix pipeline equivalent to:

\begin{verbatim}
ls | wc
\end{verbatim}

using \texttt{pipe()}, \texttt{fork()}, \texttt{dup2()} and \texttt{exec()}.

\section*{Program}

\begin{lstlisting}
#include <stdio.h>
#include <stdlib.h>
#include <unistd.h>
#include <sys/wait.h>

int main(void)
{
    int fd[2];
    pid_t pid1, pid2;

    if (pipe(fd) == -1)
    {
        perror("pipe");
        return EXIT_FAILURE;
    }

    pid1 = fork();

    if (pid1 == -1)
    {
        perror("fork");
        return EXIT_FAILURE;
    }

    if (pid1 == 0)
    {
        dup2(fd[1], STDOUT_FILENO);

        close(fd[0]);
        close(fd[1]);

        execl("/bin/ls",
              "ls",
              NULL);

        perror("execl ls");
        _exit(EXIT_FAILURE);
    }

    pid2 = fork();

    if (pid2 == -1)
    {
        perror("fork");
        return EXIT_FAILURE;
    }

    if (pid2 == 0)
    {
        dup2(fd[0], STDIN_FILENO);

        close(fd[0]);
        close(fd[1]);

        execl("/usr/bin/wc",
              "wc",
              NULL);

        perror("execl wc");
        _exit(EXIT_FAILURE);
    }

    close(fd[0]);
    close(fd[1]);

    waitpid(pid1, NULL, 0);
    waitpid(pid2, NULL, 0);

    return EXIT_SUCCESS;
}
\end{lstlisting}

% =========================================================
% COMMON VIVA
% =========================================================

\chapter*{Comprehensive Viva-Voce Questions}
\addcontentsline{toc}{chapter}{Comprehensive Viva-Voce Questions}

\begin{enumerate}

    \item What is a process?

    \item Differentiate between a program and a process.

    \item What is a Process Control Block?

    \item What is a PID?

    \item What is a PPID?

    \item Explain the \texttt{fork()} system call.

    \item Why does \texttt{fork()} return twice?

    \item What is Copy-on-Write?

    \item What is the purpose of \texttt{wait()}?

    \item Differentiate between \texttt{wait()} and \texttt{waitpid()}.

    \item What is a zombie process?

    \item What is an orphan process?

    \item What is Inter-Process Communication?

    \item Differentiate between pipe and FIFO.

    \item What are the two file descriptors returned by \texttt{pipe()}?

    \item What is shared memory?

    \item Why is shared memory generally fast?

    \item What is a message queue?

    \item Explain \texttt{msgsnd()} and \texttt{msgrcv()}.

    \item What is \texttt{exec()}?

    \item Does \texttt{exec()} create a new process?

    \item What happens to PID after a successful \texttt{exec()}?

    \item What is CPU scheduling?

    \item Differentiate preemptive and non-preemptive scheduling.

    \item Explain FCFS scheduling.

    \item What is the convoy effect?

    \item Explain SJF scheduling.

    \item Why is SJF considered optimal?

    \item What is starvation?

    \item How can aging prevent starvation?

    \item What is virtual memory?

    \item What is a page fault?

    \item Explain FIFO page replacement.

    \item What is Belady's anomaly?

    \item Explain LRU page replacement.

    \item Differentiate FIFO and LRU page replacement.

\end{enumerate}

% =========================================================
% LAB EVALUATION SHEET
% =========================================================

\chapter*{Laboratory Evaluation Record}
\addcontentsline{toc}{chapter}{Laboratory Evaluation Record}

\begin{longtable}{|c|p{6cm}|c|c|c|c|}
\hline

\textbf{Week} &
\textbf{Experiment} &
\textbf{Code} &
\textbf{Execution} &
\textbf{Viva} &
\textbf{Signature}
\\
\hline

1 & fork() & & & & \\[0.7cm]\hline
2 & Parent/Child Computation & & & & \\[0.7cm]\hline
3 & wait() & & & & \\[0.7cm]\hline
4 & Orphan/Zombie & & & & \\[0.7cm]\hline
5 & Pipe & & & & \\[0.7cm]\hline
6 & FIFO & & & & \\[0.7cm]\hline
7 & Message Queue & & & & \\[0.7cm]\hline
8 & Shared Memory & & & & \\[0.7cm]\hline
9 & FCFS & & & & \\[0.7cm]\hline
10 & SJF & & & & \\[0.7cm]\hline
11 & Priority Scheduling & & & & \\[0.7cm]\hline
12 & FIFO Page Replacement & & & & \\[0.7cm]\hline
13 & LRU Page Replacement & & & & \\[0.7cm]\hline

\end{longtable}

% =========================================================
% REFERENCES
% =========================================================

\chapter*{References}
\addcontentsline{toc}{chapter}{References}

\begin{enumerate}

    \item Abraham Silberschatz, Peter B. Galvin and Greg Gagne,
    \textit{Operating System Concepts}, Wiley.

    \item Andrew S. Tanenbaum and Herbert Bos,
    \textit{Modern Operating Systems}, Pearson.

    \item W. Richard Stevens and Stephen A. Rago,
    \textit{Advanced Programming in the UNIX Environment},
    Addison-Wesley.

    \item Michael Kerrisk,
    \textit{The Linux Programming Interface},
    No Starch Press.

    \item Linux Programmer's Manual:
    \texttt{fork(2)}, \texttt{wait(2)}, \texttt{pipe(2)},
    \texttt{execve(2)}, \texttt{msgget(2)}, and
    \texttt{shmget(2)}.

\end{enumerate}

% =========================================================
% END
% =========================================================

\end{document}